\documentclass[border=10pt]{standalone}
\usepackage{fontspec}
\AtBeginDocument{\hyphenation{zenodo}}% never hyphenate "zenodo" inside a DOI
\setmainfont{Noto Serif}[Ligatures=TeXOff]
\XeTeXgenerateactualtext=1
\newfontfamily\bengalifont{Noto Serif Bengali}[Script=Bengali, NFSSFamily=bnfam]
\usepackage[Latin,Bengali]{ucharclasses}
% group-free font switching (as in header-common.tex): robust inside braces, cells and quotes
\makeatletter
\def\BN@saved{\rmdefault}
\def\BN@bn{bnfam}
\newcommand\BNin{\edef\BN@tmp{\f@family}\ifx\BN@tmp\BN@bn\else\let\BN@saved\BN@tmp\fi\fontfamily{bnfam}\selectfont}
\newcommand\BNout{\fontfamily{\BN@saved}\selectfont}
\makeatother
\setTransitionsFor{Bengali}{\BNin}{\BNout}
% Typographic quotes and dashes are not in a ucharclasses class (class 0), so a Bangla word right
% after “ or ‘ would not trigger the Bangla font. Treat them as Latin so that every switch fires.
\XeTeXcharclass"2018=\XeTeXcharclass`A
\XeTeXcharclass"2019=\XeTeXcharclass`A
\XeTeXcharclass"201C=\XeTeXcharclass`A
\XeTeXcharclass"201D=\XeTeXcharclass`A
\XeTeXcharclass"2013=\XeTeXcharclass`A
\XeTeXcharclass"2014=\XeTeXcharclass`A
\XeTeXcharclass"2026=\XeTeXcharclass`A
\usepackage{newunicodechar}
\newunicodechar{।}{{\bengalifont\char"0964}}
\newunicodechar{॥}{{\bengalifont\char"0965}}

\usepackage[table]{xcolor}
\usepackage{booktabs,tabularx,array}
\definecolor{bapilrule}{HTML}{1F2937}
\definecolor{bapilstripe}{HTML}{F3F4F6}
\renewcommand{\arraystretch}{1.28}
\setlength{\tabcolsep}{5pt}

\begin{document}
\begin{minipage}{22cm}
{\large\bfseries সারণি ৪৫।}\ {\large কৌশলগত দিগন্ত, ২০২৬ থেকে ২০৪১ ও তার পরে}\par\vspace{7pt}
\small
\rowcolors{2}{bapilstripe}{white}
\begin{tabularx}{\linewidth}{@{}>{\hsize=0.675\hsize\raggedright\arraybackslash}X>{\hsize=0.525\hsize\raggedright\arraybackslash}X>{\hsize=1.100\hsize\raggedright\arraybackslash}X>{\hsize=1.700\hsize\raggedright\arraybackslash}X>{\hsize=0.950\hsize\raggedright\arraybackslash}X>{\hsize=1.050\hsize\raggedright\arraybackslash}X@{}}
\arrayrulecolor{bapilrule}\toprule
\rowcolor{white}\textbf{দিগন্ত} & \textbf{সময়কাল} & \textbf{কৌশলগত লক্ষ্য} & \textbf{অগ্রাধিকারমূলক পদক্ষেপ} & \textbf{প্রাতিষ্ঠানিক রূপ} & \textbf{পরবর্তী ধাপে যাওয়ার সংকেত} \\
\midrule
\textbf{H0} ভিত্তি & চতুর্থ ত্রৈমাসিক ২০২৬ & নিরাপদ ও উন্মুক্ত সূচনার প্রস্তুতি & এই ব্লুপ্রিন্ট ও কোড প্রকাশ; আইসিটি বিভাগ ও ডেটা-রক্ষকদের অবহিতকরণ; জাতীয় এআই নীতি প্রণয়নে উপকরণ প্রদান; ধাপ ০-এর অর্থায়ন & স্বাধীন গবেষণা উদ্যোগ & সিদ্ধান্ত D1 ও D5 গৃহীত \\
\textbf{H1} স্বল্পমেয়াদি & ২০২৭ থেকে ২০২৮ & উন্মুক্ত ডেটা দিয়ে মূল্য প্রমাণ & ধাপ ০ (২০২৭-এর প্রথমার্ধ): স্বাগতিক নির্বাচন, আইনি পর্যালোচনা ও ব্যয়-নিরূপণ; তারপর সনদ, ERRC ও নিবন্ধন; P1 ও P2 পাইলট; BanglaPolicyQA v0; ২০০ কর্মকর্তা প্রশিক্ষিত; মধ্যমেয়াদি মূল্যায়ন & সমঝোতা স্মারকের অধীনে বিশ্ববিদ্যালয়-ভিত্তিক ল্যাব (বিকল্প গ) & P1 ও P2 G3 উত্তীর্ণ; ইতিবাচক মধ্যমেয়াদি মূল্যায়ন \\
\textbf{H2} মধ্যমেয়াদি & ২০২৯ থেকে ২০৩০ & গ্রহণ, সম্প্রসারণ ও জাতীয় নীতির সঙ্গে সামঞ্জস্য & সংস্থাগুলোর দ্বারা P1 ও P2-এর G4 গ্রহণ; যৌক্তিক হলে BPFM পর্যায় ১; বেঞ্চমার্ক v1; ব্লুপ্রিন্ট ২.০; ২০৩০-এ সমাপ্য জাতীয় এআই নীতিচক্রের সঙ্গে সামঞ্জস্য & বহুবর্ষী মূল অর্থায়নসহ ল্যাব; সংস্থার সঙ্গে সেবা-চুক্তি & সরঞ্জাম পরিচালনাগত ব্যবহারে; টেকসই অর্থায়ন সম্মত \\
\textbf{H3} দীর্ঘমেয়াদি & ২০৩১ থেকে ২০৩৫ & জাতীয় নীতি-বুদ্ধিমত্তা নেটওয়ার্ক গঠন & মন্ত্রণালয় ও বিভাগজুড়ে সমন্বিত অনুসন্ধান; জনসম্পদ হিসেবে রক্ষিত বাংলা নীতি মডেল; সরকারজুড়ে প্রমিত আশ্বাস-পদ্ধতি; আঞ্চলিক প্রশিক্ষণ কেন্দ্র & স্বাধীন জনস্বার্থ ফাউন্ডেশন (বিকল্প ঘ) বা সংবিধিবদ্ধ কেন্দ্র বিবেচনা & BAPIL-এর বাইরেও আশ্বাস-পদ্ধতির ব্যবহার; অধিকাংশ মন্ত্রণালয়ে নেটওয়ার্ক \\
\textbf{H4} রূপকল্প & ২০৩৬ থেকে ২০৪১ & জ্ঞানভিত্তিক অর্থনীতির জন্য আস্থাভাজন এআই-সমর্থিত নীতি-ব্যবস্থা & নীতিচক্রে প্রমাণ-সেবার অন্তর্ভুক্তি; বিভিন্ন ক্ষেত্রে উন্মুক্ত বাংলা ভাষা-সম্পদ; আঞ্চলিক ও দক্ষিণ-দক্ষিণ সহযোগিতা; বাংলাদেশের দীর্ঘমেয়াদি উন্নয়ন-লক্ষ্য ও SDG-তে অবদান & আন্তর্জাতিক অংশীদারত্বসহ পরিণত জাতীয় প্রতিষ্ঠান & ২০৪১-এর স্বাধীন পর্যালোচনায় জনমূল্য নিশ্চিত \\
\textbf{২০৪১-এর পরে} & ২০৪১-এর পর & সক্ষমতাকে অভিযোজনক্ষম ও জবাবদিহিমূলক রাখা & মিশন ও সনদের পর্যায়ক্রমিক নবায়ন; কাজটি সরকারে পুরোপুরি আত্তীকৃত হলে সমাপ্তি বা একীভবন; উন্মুক্ত সংরক্ষণাগার রক্ষা & ২০৪১-এর পর্যালোচনায় নির্ধারিত & সনদ নবায়নের সিদ্ধান্ত \\
\bottomrule
\end{tabularx}
\par\vspace{5pt}{\footnotesize\itshape দিগন্ত H1 ও H2 অধ্যায় ১৫-এর ৩৬ মাসের রোডম্যাপের সঙ্গে সংগতিপূর্ণ (২০২৭-এর প্রথমার্ধে ধাপ ০; নির্দেশক মাস ১ জুলাই ২০২৭)। পরবর্তী দিগন্তগুলো কৌশলগত অভিপ্রায়, যা সারণি ৪৮-এর পর্যালোচনা-বিন্দুতে পুনর্বিবেচিত হবে।}
\par\vspace{3pt}{\scriptsize\color{gray} উৎস: BAPIL টেকনিক্যাল ব্লুপ্রিন্ট, সংস্করণ ১.০ (২০২৬)। CC BY 4.0।}
\end{minipage}
\end{document}
